\ifarxivmode
  \beginappendix
\else
  \appendices
\fi

\section{Auto-Research Tool Inventory}
\label{sec:appendix_inventory}

This appendix provides a comprehensive inventory of all surveyed works, organized by stage.

\subsection{Phase 1: Creation}
% ==================== Appendix Table: S1 Idea Generation ====================
\begingroup
\renewcommand{\arraystretch}{1.15}
\setlength{\tabcolsep}{3pt}
\footnotesize
\rowcolors{2}{white}{S1color!6}
\begin{table*}[!ht]
\centering
\caption{\textbf{Comprehensive inventory: \Sone Idea Generation.} $^\dagger$Evaluation information uncertain.}
\vspace{-7pt}
\label{tab:appendix_s1}
\stagecard{\Sone}{Idea Generation}{S1color}{figures/icons/s1_ideation.png}{%
Generating, refining, and evaluating novel research hypotheses. Techniques range from knowledge graph reasoning and retrieval-augmented generation to Multi-Agent collaboration for structured hypothesis formation.}
\vspace{2pt}
\resizebox{\linewidth}{!}{\begin{tabular}{c| >{\centering\arraybackslash}p{3cm} r r |c| >{\centering\arraybackslash}p{2.4cm} |c| p{7.8cm}}
\toprule
\rowcolor{tableheader!10}
\textbf{\#} & \textbf{Method} & \textbf{Ref} & \textbf{Venue} & \textbf{Link} & \textbf{Category} & \textbf{GitHub} & \textbf{Evaluation} \\
\midrule
\multicolumn{8}{@{}l}{\cellcolor{S1color!65}\textbf{\textsf{\textcolor{white}{~~LLM Internal Knowledge-Based Generation}}}} \\
\addlinespace[1pt]
{1} & Chain of Ideas & \cite{li2024chainofideas} & {\small arXiv'24}  & \href{https://arxiv.org/abs/2410.13185}{\faExternalLink} & LLM Internal    & \githubicon{https://github.com/DAMO-NLP-SG/CoI-Agent}                                   & Comparable to human quality; \$0.50/idea min cost \\
{2} & ResearchAgent & \cite{baek2024researchagent} & {\small NAACL'25} & \href{https://aclanthology.org/2025.naacl-long.342/}{\faExternalLink} & LLM Internal    & \githubicon{https://github.com/JinheonBaek/ResearchAgent}                                & Human + model eval; academic graph feedback \\
{3} & SciMON & \cite{wang2024scimon} & {\small ACL'24} & \href{https://arxiv.org/abs/2305.14259}{\faExternalLink} & LLM Internal    & \githubicon{https://github.com/EagleW/CLBD}                                             & Mitigates shallow novelty; iterative refinement \\
{4} & Idea Gen Agent & \cite{si2024ideas} & {\small arXiv'24} & \href{https://arxiv.org/abs/2409.04109}{\faExternalLink} & LLM Internal    & -                                                                                     & 100+ NLP researchers; LLM ideas higher novelty ($p<0.05$) \\
{5} & IRIS & \cite{iris2025} & {\small ACL'25} & \href{https://aclanthology.org/2025.acl-demo.57/}{\faExternalLink} & LLM Internal    & \githubicon{https://github.com/Anikethh/IRIS-Interactive-Research-Ideation-System}       & MCTS adaptive reasoning; human-in-the-loop platform \\
{6} & Spark & \cite{sanyal2025spark} & {\small ICCC'25} & \href{https://arxiv.org/abs/2504.20090}{\faExternalLink} & LLM Internal    & -                                                                                     & Judge model trained on 600K OpenReview reviews \\
\addlinespace[3pt]
\multicolumn{8}{@{}l}{\cellcolor{S1color!65}\textbf{\textsf{\textcolor{white}{~~External Signal-Driven Generation}}}} \\
\addlinespace[1pt]
{7} & MOOSE-Chem & \cite{yang2024moosechem} & {\small ICLR'25} & \href{https://openreview.net/forum?id=X9OfMNNepI}{\faExternalLink} & External Signal  & -                                                                                     & Rediscovers hypotheses from 51 high-impact papers \\
{8} & Nova & \cite{hu2024nova} & {\small arXiv'24} & \href{https://arxiv.org/abs/2410.14255}{\faExternalLink} & External Signal  & -                                                                                     & 3.4$\times$ more novel ideas; 2.5$\times$ more top-rated \\
{9} & SciAgents & \cite{ghafarollahi2024sciagents} & {\small arXiv'24} & \href{https://arxiv.org/abs/2409.05556}{\faExternalLink} & External Signal  & \githubicon{https://github.com/lamm-mit/SciAgentsDiscovery}                              & Multi-agent reasoning over knowledge graphs \\
{10} & SciPIP & \cite{wang2024scipip} & {\small arXiv'24} & \href{https://arxiv.org/abs/2410.23166}{\faExternalLink} & External Signal  & \githubicon{https://github.com/cheerss/SciPIP}                                           & Multi-domain; paper-anchored idea generation \\
{11} & IdeaSynth & \cite{pu2025ideasynth} & {\small CHI'25} & \href{https://arxiv.org/abs/2410.04025}{\faExternalLink} & External Signal  & -                                                                                     & 20-user study; more alternatives explored vs baseline \\
{12} & MOOSE-Chem2 & \cite{yang2025moosechem2} & {\small NeurIPS'25} & \href{https://nips.cc/virtual/2025/poster/118171}{\faExternalLink} & External Signal  & -                                                                                     & Fine-grained, experimentally actionable hypotheses \\
\addlinespace[3pt]
\multicolumn{8}{@{}l}{\cellcolor{S1color!65}\textbf{\textsf{\textcolor{white}{~~Multi-Agent Collaborative Generation}}}} \\
\addlinespace[1pt]
{13} & Combi. Creativity & \cite{gu2024combinatorial} & {\small arXiv'24} & \href{https://arxiv.org/abs/2412.14141}{\faExternalLink} & Multi-Agent   & -                                                                                     & +7--10\% similarity scores; cross-domain composition \\
{14} & Deep Ideation & \cite{zhao2025deepideation} & {\small arXiv'25} & \href{https://arxiv.org/abs/2511.02238}{\faExternalLink} & Multi-Agent      & \githubicon{https://github.com/kyZhao-1/Deep-Ideation}                                   & +10.67\% quality; surpasses conference acceptance levels \\
{15} & VirSci & \cite{su2024virsci} & {\small ACL'25} & \href{https://aclanthology.org/2025.acl-long.1368/}{\faExternalLink} & Multi-Agent      & \githubicon{https://github.com/open-sciencelab/Virtual-Scientists}                       & Outperforms single-agent on novelty$^\dagger$ \\
{16} & Multi-Agent Dial. & \cite{sigdial2025multiagent} & {\small SIGDIAL'25} & \href{https://arxiv.org/abs/2507.08350}{\faExternalLink} & Multi-Agent   & -                                                                                     & Optimal at 3 critique-revision rounds$^\dagger$ \\
{17} & Artificial Hivemind & \cite{jiang2025artificialhivemind} & {\small NeurIPS'25} & \href{https://arxiv.org/abs/2510.22954}{\faExternalLink} & Multi-Agent & -                                                                                     & 26K queries; diversity collapse across models \\
\addlinespace[3pt]
\multicolumn{8}{@{}l}{\cellcolor{S1color!65}\textbf{\textsf{\textcolor{white}{~~Novelty and Feasibility Assessment}}}} \\
\addlinespace[1pt]
{18} & IdeaBench & \cite{guo2025ideabench} & {\small KDD'25} & \href{https://doi.org/10.1145/3711896.3737419}{\faExternalLink} & Evaluation       & -                                                                                     & 2,374 papers; 8 domains; novelty $>$0.6, feasibility $<$0.5 \\
{19} & LiveIdeaBench & \cite{liveideabench2024} & {\small arXiv'24} & \href{https://arxiv.org/abs/2412.17596}{\faExternalLink} & Evaluation       & -                                                                                     & 40+ models; 1,180 keywords; 22 scientific domains \\
{20} & AI Idea Bench 2025 & \cite{aiideabench2025} & {\small arXiv'25} & \href{https://arxiv.org/abs/2504.14191}{\faExternalLink} & Evaluation       & \githubicon{https://github.com/yansheng-qiu/AI_Idea_Bench_2025}                          & 3,495 papers; alignment + general reference eval \\
{21} & HeurekaBench & \cite{heurekabench2026} & {\small ICLR'26} & \href{https://arxiv.org/abs/2601.01678}{\faExternalLink} & Evaluation       & \githubicon{https://github.com/mlbio-epfl/HeurekaBench}                                  & +22\% with critic module; open-ended science tasks \\
{22} & ResearchBench & \cite{researchbench2025} & {\small ACL'26} & \href{https://arxiv.org/abs/2503.21248}{\faExternalLink} & Evaluation       & -                                                                                     & 12 disciplines; inspiration retrieval + ranking \\
{23} & HindSight & \cite{hindsight2026} & {\small arXiv'26} & \href{https://arxiv.org/abs/2603.15164}{\faExternalLink} & Evaluation       & -                                                                                     & LLM novelty negatively correlated with impact ($\rho$=$-$0.29) \\
{24} & Rubric Rewards & \cite{rubricrewards2025} & {\small arXiv'25} & \href{https://arxiv.org/abs/2512.23707}{\faExternalLink} & LLM Internal & - & 70\% expert preference; RL with rubric self-grading \\
{25} & DeepInnovator & \cite{deepinnovator2026} & {\small arXiv'26} & \href{https://arxiv.org/abs/2602.18920}{\faExternalLink} & LLM Internal & \githubicon{https://github.com/HKUDS/DeepInnovator} & 80--94\% win rates vs frontier; 14B model \\
{26} & FlowPIE & \cite{flowpie2026} & {\small arXiv'26} & \href{https://arxiv.org/abs/2603.29557}{\faExternalLink} & External Signal & - & Higher novelty, feasibility, diversity vs baselines \\
\bottomrule
\end{tabular}}
\vspace{-1cm}
\end{table*}
\endgroup


\clearpage
% ==================== Appendix Table: S2 Literature Review ====================
\begingroup
\renewcommand{\arraystretch}{1.15}
\setlength{\tabcolsep}{3pt}
\footnotesize
\rowcolors{2}{white}{S2color!6}
\begin{table*}[!ht]
\centering
\caption{\textbf{Comprehensive inventory: \Stwo Literature Review.} $^\dagger$Evaluation information uncertain.}
\vspace{-7pt}
\label{tab:appendix_s2}
\stagecard{\Stwo}{Literature Review}{S2color}{figures/icons/s2_literature.png}{%
Retrieving, synthesizing, and organizing prior work into coherent narratives. Modern approaches span semantic retrieval, citation-graph traversal, and Deep Research agents that iteratively explore the literature.}
\vspace{2pt}
\resizebox{\linewidth}{!}{\begin{tabular}{c| >{\centering\arraybackslash}p{3cm} r r |c| >{\centering\arraybackslash}p{2.3cm} |c| p{8cm}}
\toprule
\rowcolor{tableheader!10}
\textbf{\#} & \textbf{Method} & \textbf{Ref} & \textbf{Venue} & \textbf{Link} & \textbf{Category} & \textbf{GitHub} & \textbf{Evaluation} \\
\midrule
\multicolumn{8}{@{}l}{\cellcolor{S2color!65}\textbf{\textsf{\textcolor{white}{~~Literature Retrieval}}}} \\
\addlinespace[1pt]
{1} & CiteME & \cite{citeme2024} & {\small arXiv'24} & \href{https://arxiv.org/abs/2407.12861}{\faExternalLink} & Retrieval        & -                                                                                     & Citation fidelity benchmark \\
{2} & LitLLM & \cite{agarwal2024litllm} & {\small arXiv'24} & \href{https://arxiv.org/abs/2402.01788}{\faExternalLink} & Retrieval        & -                                                                                     & LLM + academic database integration \\
{3} & LitSearch & \cite{litsearch2024} & {\small arXiv'24} & \href{https://arxiv.org/abs/2407.18940}{\faExternalLink} & Retrieval        & \githubicon{https://github.com/princeton-nlp/LitSearch}                                  & Retrieval precision benchmark \\
{4} & PaperQA2 & \cite{skarlinski2024paperqa2} & {\small arXiv'24} & \href{https://arxiv.org/abs/2409.13740}{\faExternalLink} & Retrieval        & \githubicon{https://github.com/Future-House/paper-qa}                                    & Matches/exceeds expert on 3 tasks; 70\% contradiction validation \\
{5} & OpenResearcher & \cite{li2024openresearcher} & {\small EMNLP'24} & \href{https://arxiv.org/abs/2408.09578}{\faExternalLink} & Retrieval        & -                                                                                     & RAG + graph traversal for literature exploration \\
{6} & PaSa & \cite{pasa2025} & {\small arXiv'25} & \href{https://arxiv.org/abs/2501.10120}{\faExternalLink} & Retrieval        & \githubicon{https://github.com/bytedance/pasa}                                           & Agentic multi-step iterative retrieval \\
\addlinespace[3pt]
\multicolumn{8}{@{}l}{\cellcolor{S2color!65}\textbf{\textsf{\textcolor{white}{~~Survey \& Related Work Generation}}}} \\
\addlinespace[1pt]
{7} & ChatPaper & \cite{chatpaper2023} & {\small GitHub'23} & \href{https://github.com/kaixindelele/ChatPaper}{\faExternalLink} & Generation       & \githubicon{https://github.com/kaixindelele/ChatPaper}                                   & 19K+ GitHub stars; arXiv summarization tool \\
{8} & PaperQA & \cite{paperqa2024github} & {\small arXiv'23} & \href{https://arxiv.org/abs/2312.07559}{\faExternalLink} & Generation       & \githubicon{https://github.com/Future-House/paper-qa}                                    & 8K+ GitHub stars; RAG for scientific Q\&A \\
{9} & AutoSurvey & \cite{wang2024autosurvey} & {\small arXiv'24} & \href{https://arxiv.org/abs/2406.10252}{\faExternalLink} & Generation       & \githubicon{https://github.com/AutoSurveys/AutoSurvey}                                   & First end-to-end LLM survey drafting system \\
{10} & GPT Researcher & \cite{gptresearcher2024} & {\small GitHub'24} & \href{https://github.com/assafelovic/gpt-researcher}{\faExternalLink} & Generation       & \githubicon{https://github.com/assafelovic/gpt-researcher}                               & 26K+ GitHub stars; comprehensive report generation \\
{11} & LLMs for Lit.\ Review & \cite{emnlp2025litreview} & {\small arXiv'24} & \href{https://arxiv.org/abs/2412.13612}{\faExternalLink} & Generation       & -                                                                                     & Hallucination analysis; models still generate errors$^\dagger$ \\
{12} & STORM & \cite{shao2024storm} & {\small arXiv'24} & \href{https://arxiv.org/abs/2402.14207}{\faExternalLink} & Generation       & \githubicon{https://github.com/stanford-oval/storm}                                      & Multi-perspective question-asking for outlines \\
{13} & Agentic AutoSurvey & \cite{agenticautosurvey2025} & {\small arXiv'25} & \href{https://arxiv.org/abs/2509.18661}{\faExternalLink} & Generation       & -                                                                                     & Multi-agent role decomposition$^\dagger$ \\
{14} & Citegeist & \cite{beger2025citegeist} & {\small arXiv'25} & \href{https://arxiv.org/abs/2503.23229}{\faExternalLink} & Generation       & -                                                                                     & Dynamic RAG pipeline on arXiv corpus \\
{15} & IterSurvey & \cite{itersurvey2025} & {\small arXiv'25} & \href{https://arxiv.org/abs/2510.21900}{\faExternalLink} & Generation       & \githubicon{https://github.com/HancCui/IterSurvey_Autosurveyv2}                          & Iterative outline planning with stability checks \\
{16} & LiRA & \cite{lira2025} & {\small arXiv'25} & \href{https://arxiv.org/abs/2510.05138}{\faExternalLink} & Generation       & -                                                                                     & Multi-agent retrieval + verification + narrative \\
{17} & SurveyForge & \cite{gao2025surveyforge} & {\small arXiv'25} & \href{https://arxiv.org/abs/2503.04629}{\faExternalLink} & Generation       & \githubicon{https://github.com/Alpha-Innovator/SurveyForge}                              & Outperforms AutoSurvey on outline quality$^\dagger$ \\
{18} & SurveyG & \cite{surveyg2025} & {\small arXiv'25} & \href{https://arxiv.org/abs/2510.07733}{\faExternalLink} & Generation       & -                                                                                     & Three-layer citation graph (Foundation/Dev/Frontier) \\
{19} & SurveyX & \cite{liang2025surveyx} & {\small arXiv'25} & \href{https://arxiv.org/abs/2502.14776}{\faExternalLink} & Generation       & -                                                                                     & +0.259 content quality improvement; near expert level \\
{20} & InteractiveSurvey & \cite{interactivesurvey2025} & {\small arXiv'25} & \href{https://arxiv.org/abs/2504.08762}{\faExternalLink} & Generation       & \githubicon{https://github.com/TechnicolorGUO/InteractiveSurvey}                        & User-customizable reference categorization + outlines \\
{21} & CiteLLM & \cite{citellm2026} & {\small arXiv'26} & \href{https://arxiv.org/abs/2602.23075}{\faExternalLink} & Generation       & -                                                                                     & Hallucination-free via trusted repository routing \\
\addlinespace[3pt]
\multicolumn{8}{@{}l}{\cellcolor{S2color!65}\textbf{\textsf{\textcolor{white}{~~Deep Research Agents}}}} \\
\addlinespace[1pt]
{22} & ASReview & \cite{asreview2020} & {\small Nature MI'21} & \href{https://www.nature.com/articles/s42256-020-00287-7}{\faExternalLink} & Deep Research & \githubicon{https://github.com/asreview/asreview}                                        & Active learning; up to 95\% effort reduction \\
{23} & CHIME & \cite{kang2024chime} & {\small arXiv'24} & \href{https://arxiv.org/abs/2407.16148}{\faExternalLink} & Deep Research    & -                                                                                     & Hierarchical organization of scientific studies \\
{24} & DeepResearch-Agent & \cite{deepresearchagent2025} & {\small GitHub'25} & \href{https://github.com/SkyworkAI/DeepResearchAgent}{\faExternalLink} & Deep Research    & \githubicon{https://github.com/SkyworkAI/DeepResearchAgent}                              & Hierarchical multi-agent; planner + sub-agents \\
{25} & DeerFlow & \cite{deerflow2025} & {\small GitHub'25} & \href{https://github.com/bytedance/deer-flow}{\faExternalLink} & Deep Research    & \githubicon{https://github.com/bytedance/deer-flow}                                      & Sub-agents with shared memory; sandboxed execution \\
{26} & OpenScholar & \cite{openscholar2025} & {\small Nature'26} & \href{https://doi.org/10.1038/s41586-025-10072-4}{\faExternalLink} & Deep Research    & -                                                                                     & 45M papers; +6.1\% over GPT-4o, +5.5\% over PaperQA2 \\
{27} & AutoAgent & \cite{tang2025autodeepresearch} & {\small arXiv'25} & \href{https://arxiv.org/abs/2502.05957}{\faExternalLink} & Deep Research    & -                                                                                     & Universal LLM compatibility; GAIA benchmark \\
{28} & Tongyi DeepResearch & \cite{tongyi2025deepresearch} & {\small GitHub'25} & \href{https://github.com/Alibaba-NLP/DeepResearch}{\faExternalLink} & Deep Research    & \githubicon{https://github.com/Alibaba-NLP/DeepResearch}                                 & 30.5B params (3.3B activated); SOTA on Deep Research \\
{29} & O-Researcher & \cite{oresearcher2026} & {\small arXiv'26} & \href{https://arxiv.org/abs/2601.03743}{\faExternalLink} & Deep Research    & -                                                                                     & Multi-agent distillation + agentic RL \\
{30} & OpenResearcher & \cite{li2026openresearcher} & {\small arXiv'26} & \href{https://arxiv.org/abs/2603.20278}{\faExternalLink} & Deep Research    & \githubicon{https://github.com/TIGER-AI-Lab/OpenResearcher}                              & 54.8\% BrowseComp-Plus; 97K+ trajectories \\
\addlinespace[3pt]
\multicolumn{8}{@{}l}{\cellcolor{S2color!65}\textbf{\textsf{\textcolor{white}{~~Retrieval and Synthesis Quality Assessment}}}} \\
\addlinespace[1pt]
{31} & DeepScholar-Bench & \cite{deepscholar2025} & {\small arXiv'25} & \href{https://arxiv.org/abs/2508.20033}{\faExternalLink} & Evaluation       & \githubicon{https://github.com/guestrin-lab/deepscholar-bench}                            & Coverage, coherence, factual accuracy benchmark \\
{32} & ReportBench & \cite{reportbench2025} & {\small arXiv'25} & \href{https://arxiv.org/abs/2508.15804}{\faExternalLink} & Evaluation       & \githubicon{https://github.com/ByteDance-BandAI/ReportBench}                             & 100-prompt benchmark from 678 filtered survey papers \\
{33} & IDRBench & \cite{idrbench2026} & {\small arXiv'26} & \href{https://arxiv.org/abs/2601.06676}{\faExternalLink} & Evaluation       & -                                                                                     & 100 tasks; interactive Deep Research evaluation \\
{34} & ScholarGym & \cite{scholargym2026} & {\small arXiv'26} & \href{https://arxiv.org/abs/2601.21654}{\faExternalLink} & Evaluation       & -                                                                                     & 2,536 queries; query planning + tool invocation \\
{35} & SciNetBench & \cite{scinetbench2026} & {\small arXiv'26} & \href{https://arxiv.org/abs/2601.03260}{\faExternalLink} & Evaluation       & -                                                                                     & 18M papers; relation-aware retrieval $<$20\% \\
\bottomrule
\end{tabular}}
\end{table*}
\endgroup


\clearpage
% ==================== Appendix Table: S3 Coding & Experiments ====================
\begingroup
\renewcommand{\arraystretch}{1.15}
\setlength{\tabcolsep}{3pt}
\footnotesize
\rowcolors{2}{white}{S3color!6}
\begin{table*}[!ht]
\centering
\caption{\textbf{Comprehensive inventory: \Sthree Coding \& Experiments.} $^\dagger$Evaluation information uncertain.}
\vspace{-7pt}
\label{tab:appendix_s3}
\stagecard{\Sthree}{Coding \& Experiments}{S3color}{figures/icons/s3_coding.png}{%
Translating ideas into executable code, running experiments at scale, and analyzing results. This stage spans code generation, Paper-to-Code translation, autonomous experiment orchestration, and result interpretation.}
\vspace{2pt}
\resizebox{\linewidth}{!}{\begin{tabular}{c| >{\centering\arraybackslash}p{3cm} r r |c| >{\centering\arraybackslash}p{2.4cm} |c| p{8cm}}
\toprule
\rowcolor{tableheader!10}
\textbf{\#} & \textbf{Method} & \textbf{Ref} & \textbf{Venue} & \textbf{Link} & \textbf{Category} & \textbf{GitHub} & \textbf{Evaluation} \\
\midrule
\multicolumn{8}{@{}l}{\cellcolor{S3color!65}\textbf{\textsf{\textcolor{white}{~~Code Generation}}}} \\
\addlinespace[1pt]
{1} & SWE-bench & \cite{swebench2024} & {\small ICLR'24} & \href{https://arxiv.org/abs/2310.06770}{\faExternalLink} & Code Gen.        & \githubicon{https://github.com/princeton-nlp/SWE-bench}                                  & 2,294 real GitHub issues; Verified split (500 problems) \\
{2} & SWE-agent & \cite{yang2024sweagent} & {\small arXiv'24} & \href{https://arxiv.org/abs/2405.15793}{\faExternalLink} & Code Gen.        & \githubicon{https://github.com/princeton-nlp/SWE-agent}                                  & Agent--computer interface paradigm for coding \\
{3} & OpenHands & \cite{wang2024openhands} & {\small ICLR'25} & \href{https://arxiv.org/abs/2407.16741}{\faExternalLink} & Code Gen.        & \githubicon{https://github.com/All-Hands-AI/OpenHands}                                   & Open platform for generalist coding agents \\
{4} & SWE-bench Pro & \cite{deng2025swebenchpro} & {\small arXiv'25} & \href{https://arxiv.org/abs/2509.16941}{\faExternalLink} & Code Gen.        & -                                                                                     & 1,865 enterprise problems; best score 23\% \\
{5} & SWE-EVO & \cite{thai2025sweevo} & {\small arXiv'25} & \href{https://arxiv.org/abs/2512.18470}{\faExternalLink} & Code Gen.        & -                                                                                     & Software evolution benchmark; best score 25\% \\
\addlinespace[3pt]
\multicolumn{8}{@{}l}{\cellcolor{S3color!65}\textbf{\textsf{\textcolor{white}{~~Paper-to-Code}}}} \\
\addlinespace[1pt]
{6} & FunSearch & \cite{funsearch2024} & {\small Nature'24} & \href{https://www.nature.com/articles/s41586-023-06924-6}{\faExternalLink} & Paper-to-Code    & \githubicon{https://github.com/google-deepmind/funsearch}                                & New cap-set solutions; evolutionary program search \\
{7} & SciCode & \cite{scicode2024} & {\small arXiv'24} & \href{https://arxiv.org/abs/2407.13168}{\faExternalLink} & Paper-to-Code    & \githubicon{https://github.com/scicode-bench/SciCode}                                    & Research-level coding across math, physics, chemistry \\
{8} & PaperBench & \cite{paperbench2025} & {\small arXiv'25} & \href{https://arxiv.org/abs/2504.01848}{\faExternalLink} & Paper-to-Code    & \githubicon{https://github.com/openai/preparedness}                                      & 20 ICML'24 papers; 8,316 gradable subtasks \\
{9} & PaperCoder & \cite{papercoder2025} & {\small arXiv'25} & \href{https://arxiv.org/abs/2504.17192}{\faExternalLink} & Paper-to-Code    & \githubicon{https://github.com/going-doer/Paper2Code}                                    & 3-stage multi-agent; ML papers to code repos \\
{10} & ResearchCodeBench & \cite{researchcodebench2025} & {\small arXiv'25} & \href{https://arxiv.org/abs/2506.02314}{\faExternalLink} & Paper-to-Code    & -                                                                                     & 212 novel ML tasks; best 37.3\% (Gemini-2.5-Pro) \\
{11} & SciReplicate-Bench & \cite{scireplicatebench2025} & {\small arXiv'25} & \href{https://arxiv.org/abs/2504.00255}{\faExternalLink} & Paper-to-Code    & \githubicon{https://github.com/xyzCS/SciReplicate-Bench}                                 & 100 tasks from 36 NLP papers; 39\% ceiling \\
\addlinespace[3pt]
\multicolumn{8}{@{}l}{\cellcolor{S3color!65}\textbf{\textsf{\textcolor{white}{~~Experiment Execution \& Orchestration}}}} \\
\addlinespace[1pt]
{12} & BioPlanner & \cite{bioplanner2024} & {\small arXiv'23} & \href{https://arxiv.org/abs/2310.10632}{\faExternalLink} & Execution        & \githubicon{https://github.com/bioplanner/bioplanner}                                    & Biological protocol planning evaluation \\
{13} & CRISPR-GPT & \cite{crisprgpt2024} & {\small arXiv'24} & \href{https://arxiv.org/abs/2404.18021}{\faExternalLink} & Execution        & -                                                                                     & Gene-editing experiment design assistance \\
{14} & DS-Agent & \cite{dsagent2024} & {\small arXiv'24} & \href{https://arxiv.org/abs/2402.17453}{\faExternalLink} & Execution        & \githubicon{https://github.com/guosyjlu/DS-Agent}                                        & End-to-end data science workflow automation \\
{15} & MLE-Bench & \cite{chan2024mlebench} & {\small arXiv'24} & \href{https://arxiv.org/abs/2410.07095}{\faExternalLink} & Execution        & -                                                                                     & 75 Kaggle competitions benchmark \\
{16} & MLAgentBench & \cite{mlagentbench2024} & {\small arXiv'24} & \href{https://arxiv.org/abs/2310.03302}{\faExternalLink} & Execution        & \githubicon{https://github.com/snap-stanford/MLAgentBench}                               & 13 ML experimentation tasks benchmark \\
{17} & MLR-Copilot & \cite{mlrcopilot2024} & {\small arXiv'24} & \href{https://arxiv.org/abs/2408.14033}{\faExternalLink} & Execution        & -                                                                                     & IdeaAgent + ExperimentAgent dual-agent pipeline \\
{18} & AIDE & \cite{jiang2025aide} & {\small arXiv'25} & \href{https://arxiv.org/abs/2502.13138}{\faExternalLink} & Execution        & -                                                                                     & SOTA on MLE-Bench + RE-Bench; tree search in code space \\
{19} & AlphaEvolve & \cite{alphaevolve2025} & {\small arXiv'25} & \href{https://arxiv.org/abs/2506.13131}{\faExternalLink} & Execution        & -                                                                                     & LLM-generated mutations + automated evaluators \\
{20} & AutoReproduce & \cite{autoreproduce2025} & {\small arXiv'25} & \href{https://arxiv.org/abs/2505.20662}{\faExternalLink} & Execution        & \githubicon{https://github.com/AI9Stars/AutoReproduce}                                   & Paper lineage algorithm for experiment reproduction \\
{21} & CURIE & \cite{curie2025} & {\small arXiv'25} & \href{https://arxiv.org/abs/2502.16069}{\faExternalLink} & Execution        & \githubicon{https://github.com/Just-Curieous/Curie}                                      & Rigorous automated experimentation framework \\
{22} & MLGym & \cite{mlgym2025} & {\small arXiv'25} & \href{https://arxiv.org/abs/2502.14499}{\faExternalLink} & Execution        & -                                                                                     & AI research agent gym benchmark \\
{23} & MLR-Bench & \cite{mlrbench2025} & {\small arXiv'25} & \href{https://arxiv.org/abs/2505.19955}{\faExternalLink} & Execution        & -                                                                                     & 201 tasks (NeurIPS/ICLR/ICML); 80\% fabrication rate \\
{24} & Execution-Grounded & \cite{si2026executiongrounded} & {\small arXiv'26} & \href{https://arxiv.org/abs/2601.14525}{\faExternalLink} & Execution       & -                                                                                     & 69.4\% vs 48.0\% GRPO; parallel GPU search \\
{25} & Learn to Discover & \cite{yuksekgonul2026learntodiscover} & {\small arXiv'26} & \href{https://arxiv.org/abs/2601.16175}{\faExternalLink} & Execution  & -                                                                                     & Test-time training + RL; math, GPU kernel, biology \\
{26} & SciNav & \cite{scinav2026} & {\small arXiv'26} & \href{https://arxiv.org/abs/2603.20256}{\faExternalLink} & Execution        & -                                                                                     & Pairwise tree-search branch selection \\
{27} & FrontierScience & \cite{wang2026frontierscience} & {\small arXiv'26} & \href{https://arxiv.org/abs/2601.21165}{\faExternalLink} & Execution        & -                                                                                     & Expert-level tasks; Olympiad + PhD difficulty \\
\addlinespace[3pt]
\multicolumn{8}{@{}l}{\cellcolor{S3color!65}\textbf{\textsf{\textcolor{white}{~~Code Correctness and Reproducibility Assessment}}}} \\
\addlinespace[1pt]
{28} & DiscoveryBench & \cite{majumder2024discoverybench} & {\small arXiv'24} & \href{https://arxiv.org/abs/2407.01725}{\faExternalLink} & Analysis         & \githubicon{https://github.com/allenai/discoverybench}                                   & Data-driven insight extraction benchmark \\
{29} & DiscoveryWorld & \cite{discoveryworld2024} & {\small arXiv'24} & \href{https://arxiv.org/abs/2406.06769}{\faExternalLink} & Analysis         & \githubicon{https://github.com/allenai/discoveryworld}                                   & 120 tasks; 8 topics; 3 difficulty levels \\
{30} & InfiAgent-DABench & \cite{infidabench2024} & {\small arXiv'24} & \href{https://arxiv.org/abs/2401.05507}{\faExternalLink} & Analysis         & -                                                                                     & End-to-end data analysis workflow benchmark \\
{31} & ScienceAgentBench & \cite{scienceagentbench2024} & {\small arXiv'24} & \href{https://arxiv.org/abs/2410.05080}{\faExternalLink} & Analysis         & -                                                                                     & Rigorous data-driven scientific discovery assessment \\
{32} & LAB-Bench & \cite{labbench2024} & {\small arXiv'24} & \href{https://arxiv.org/abs/2407.10362}{\faExternalLink} & Execution & \githubicon{https://github.com/Future-House/LAB-Bench} & Multi-domain biology research task benchmark \\
{33} & KernelBench & \cite{kernelbench2025} & {\small arXiv'25} & \href{https://arxiv.org/abs/2502.10517}{\faExternalLink} & Execution & \githubicon{https://github.com/ScalingIntelligence/KernelBench} & GPU kernel generation benchmark \\
{34} & TritonBench & \cite{tritonbench2025} & {\small arXiv'25} & \href{https://arxiv.org/abs/2502.14752}{\faExternalLink} & Execution & \githubicon{https://github.com/thunlp/TritonBench} & Triton operator generation benchmark \\
{35} & AstaBench & \cite{astabench2025} & {\small arXiv'25} & \href{https://arxiv.org/abs/2510.21652}{\faExternalLink} & Execution & \githubicon{https://github.com/allenai/asta-bench} & 2,400+ problems; multi-domain scientific research \\
{36} & ResearchClawBench & \cite{researchclawbench2025} & {\small arXiv'25} & \href{https://arxiv.org/abs/2512.16969}{\faExternalLink} & Execution & \githubicon{https://github.com/InternScience/ResearchClawBench} & Scientist-aligned workflow benchmark \\
{37} & EXP-Bench & \cite{expbench2025} & {\small ICLR'26} & \href{https://openreview.net/forum?id=KjgyAm383Z}{\faExternalLink} & Execution & \githubicon{https://github.com/Just-Curieous/Curie/tree/main/benchmark/exp_bench} & 461 tasks from 51 AI papers \\
{38} & PostTrainBench & \cite{posttrainbench2026} & {\small arXiv'26} & \href{https://arxiv.org/abs/2603.08640}{\faExternalLink} & Execution & \githubicon{https://github.com/aisa-group/PostTrainBench} & LLM post-training automation benchmark \\
\bottomrule
\end{tabular}}
\end{table*}
\endgroup


\clearpage
% ==================== Appendix Table: S4 Tables & Figures ====================
\begingroup
\renewcommand{\arraystretch}{1.15}
\setlength{\tabcolsep}{3pt}
\footnotesize
\rowcolors{2}{white}{S4color!6}
\begin{table*}[!ht]
\centering
\caption{\textbf{Comprehensive inventory: \Sfour Tables \& Figures.} $^\dagger$Evaluation information uncertain.}
\vspace{-7pt}
\label{tab:appendix_s4}
\stagecard{\Sfour}{Tables \& Figures}{S4color}{figures/icons/s4_figures.png}{%
Creating Method Diagrams, result plots, comparison tables, and LaTeX formulas. Scientific visualization transforms raw experimental outputs into publication-quality charts, illustrations, and structured tables.}
\vspace{2pt}
\resizebox{\linewidth}{!}{\begin{tabular}{c|>{\centering\arraybackslash}p{3cm} r r |c| >{\centering\arraybackslash}p{2.7cm} |c| p{6.9cm}}
\toprule
\rowcolor{tableheader!10}
\textbf{\#} & \textbf{Method} & \textbf{Ref} & \textbf{Venue} & \textbf{Link} & \textbf{Category} & \textbf{GitHub} & \textbf{Evaluation} \\
\midrule
\multicolumn{8}{@{}l}{\cellcolor{S4color!65}\textbf{\textsf{\textcolor{white}{~~Scientific Figure Generation}}}} \\
\addlinespace[1pt]
{1} & ChartGPT & \cite{yuan2023chartgpt} & {\small arXiv'23} & \href{https://arxiv.org/abs/2311.01920}{\faExternalLink} & Data Viz         & -                                                                                     & 6-step reasoning for chart generation \\
{2} & MatPlotAgent & \cite{matplotagent2024} & {\small arXiv'24} & \href{https://arxiv.org/abs/2402.11453}{\faExternalLink} & Data Viz         & -                                                                                     & +12.3 over GPT-4 base; VLM visual feedback$^\dagger$ \\
{3} & CoDA & \cite{coda2025} & {\small arXiv'25} & \href{https://arxiv.org/abs/2510.03194}{\faExternalLink} & Data Viz         & -                                                                                     & +41.5\% over baselines; multi-agent collaboration \\
{4} & PlotGen & \cite{plotgen2025} & {\small arXiv'25} & \href{https://arxiv.org/abs/2502.00988}{\faExternalLink} & Data Viz         & -                                                                                     & 4--6\% improvement over baselines$^\dagger$ \\
{5} & VIS-Shepherd & \cite{visshepherd2025} & {\small arXiv'25} & \href{https://arxiv.org/abs/2506.13326}{\faExternalLink} & Figure Editing   & -                                                                                     & Constructive critique feedback framework \\
{6} & DiagramAgent & \cite{diagramagent2024} & {\small CVPR'25} & \href{https://arxiv.org/abs/2411.11916}{\faExternalLink} & Data Viz         & -                                                                                     & 4 specialized agents; 8 diagram categories \\
{7} & StarVector & \cite{starVector2025} & {\small CVPR'25} & \href{https://arxiv.org/abs/2312.11556}{\faExternalLink} & Method Diagrams  & -                                                                                     & Scalable SVG generation from descriptions$^\dagger$ \\
{8} & VisCoder & \cite{viscoder2025} & {\small EMNLP'25} & \href{https://arxiv.org/abs/2506.03930}{\faExternalLink} & Data Viz         & -                                                                                     & VisCode-200K dataset; 90\%+ execution pass rate$^\dagger$ \\
{9} & AI-Generated Figures & \cite{aifigurepolicies2026} & {\small arXiv'26} & \href{https://arxiv.org/abs/2603.16159}{\faExternalLink} & Policy           & -                                                                                     & Publisher policy survey (Nature, Science, etc.) \\
{10} & AutoFigure-Edit & \cite{autofigure2026} & {\small arXiv'26} & \href{https://arxiv.org/abs/2603.06674}{\faExternalLink} & Method Diagrams  & \githubicon{https://github.com/ResearAI/AutoFigure-Edit}                & Editable text-to-SVG scientific illustrations$^\dagger$ \\
{11} & AutoFigure & \cite{autofigure2026iclr} & {\small ICLR'26} & \href{https://arxiv.org/abs/2602.03828}{\faExternalLink} & Method Diagrams  & \githubicon{https://github.com/ResearAI/AutoFigure}                     & FigureBench (3,300 pairs); publication-ready illust. \\
{12} & PaperBanana & \cite{paperbanana2026} & {\small arXiv'26} & \href{https://arxiv.org/abs/2601.23265}{\faExternalLink} & Method Diagrams  & -                                                                                     & 292 test cases; outperforms baselines$^\dagger$ \\
{13} & SAIL & \cite{sail2026} & {\small arXiv'26} & \href{https://arxiv.org/abs/2603.18145}{\faExternalLink} & Figure Editing   & -                                                                                     & Domain logic / code syntax separation \\
\addlinespace[3pt]
\multicolumn{8}{@{}l}{\cellcolor{S4color!65}\textbf{\textsf{\textcolor{white}{~~Table Understanding \& Generation}}}} \\
\addlinespace[1pt]
{14} & ArxivDIGESTables & \cite{arxivdigestables2024} & {\small EMNLP'24} & \href{https://arxiv.org/abs/2410.22360}{\faExternalLink} & Table Gen.        & -                                                                                     & Literature comparison table synthesis \\
{15} & Chain-of-Table & \cite{chainoftable2024} & {\small ICLR'24} & \href{https://arxiv.org/abs/2401.04398}{\faExternalLink} & Table Reasoning  & -                                                                                     & Multi-step table reasoning chains \\
{16} & ShowTable & \cite{showtable2025} & {\small CVPR'26} & \href{https://arxiv.org/abs/2512.13303}{\faExternalLink} & Table Viz        & -                                                                                     & Collaborative reflection and refinement$^\dagger$ \\
{17} & Table2LaTeX-RL & \cite{table2latexrl2025} & {\small arXiv'25} & \href{https://arxiv.org/abs/2509.17589}{\faExternalLink} & Table Conversion & -                                                                                     & Image-to-LaTeX via reinforced multimodal LM \\
\addlinespace[3pt]
\multicolumn{8}{@{}l}{\cellcolor{S4color!65}\textbf{\textsf{\textcolor{white}{~~Mathematical Formulas \& TikZ}}}} \\
\addlinespace[1pt]
{18} & AutomaTikZ & \cite{belouadi2024automatikz} & {\small ICLR'24} & \href{https://arxiv.org/abs/2310.00367}{\faExternalLink} & TikZ             & -                                                                                     & DaTikZ: first large-scale dataset (120K drawings) \\
{19} & DeTikZify & \cite{belouadi2024detikzify} & {\small NeurIPS'24} & \href{https://arxiv.org/abs/2405.15306}{\faExternalLink} & TikZ            & -                                                                                     & 360K TikZ graphics; MCTS iterative refinement \\
{20} & TikZilla & \cite{tikzilla2026} & {\small arXiv'26} & \href{https://arxiv.org/abs/2603.03072}{\faExternalLink} & TikZ             & -                                                                                     & 3B/8B matches GPT-5; SFT+RL on expanded DaTikZ \\
\addlinespace[3pt]
\multicolumn{8}{@{}l}{\cellcolor{S4color!65}\textbf{\textsf{\textcolor{white}{~~Visual Fidelity and Scientific Accuracy Assessment}}}} \\
\addlinespace[1pt]
{21} & PlotCraft & \cite{plotcraft2025} & {\small arXiv'25} & \href{https://arxiv.org/abs/2511.00010}{\faExternalLink} & Benchmark         & -                                                                                     & 1K-task benchmark; 48 chart types \\
{22} & TeXpert & \cite{texpert2025} & {\small SDP'25} & \href{https://aclanthology.org/2025.sdp-1.2/}{\faExternalLink} & Benchmark         & -                                                                                     & 3-level difficulty; 78.8\%/58.7\%/17.5\%$^\dagger$ \\
{23} & AbGen & \cite{abgen2025} & {\small ACL'25} & \href{https://arxiv.org/abs/2507.13300}{\faExternalLink} & Benchmark         & -                                                                                     & 1,500 ablation studies; 807 NLP papers \\
{24} & SciFig & \cite{scifig2026} & {\small arXiv'26} & \href{https://arxiv.org/abs/2601.04390}{\faExternalLink} & Benchmark         & -                                                                                     & Rubric-based evaluation; 2K+ pipeline figures \\
{25} & SciFlow-Bench & \cite{scifigbench2026} & {\small arXiv'26} & \href{https://arxiv.org/abs/2602.09809}{\faExternalLink} & Benchmark         & -                                                                                     & 500 framework figures; inverse-parsing evaluation \\
{26} & FigureBench & \cite{autofigure2026iclr} & {\small ICLR'26} & \href{https://arxiv.org/abs/2602.03828}{\faExternalLink} & Benchmark         & \githubicon{https://github.com/ResearAI/AutoFigure}                     & 3,300 text-figure pairs; publication-ready eval \\
\bottomrule
\end{tabular}}
\end{table*}
\endgroup



\clearpage
\subsection{Phase 2: Writing}

% ==================== Appendix Table: S5 Paper Writing ====================
\begingroup
\renewcommand{\arraystretch}{1.15}
\setlength{\tabcolsep}{3pt}
\footnotesize
\rowcolors{2}{white}{S5color!6}
\begin{table*}[!ht]
\centering
\caption{\textbf{Comprehensive inventory: \Sfive Paper Writing.} $^\dagger$Evaluation information uncertain.}
\vspace{-7pt}
\label{tab:appendix_s5}
\stagecard{\Sfive}{Paper Writing}{S5color}{figures/icons/s5_writing.png}{%
Drafting, editing, and polishing academic manuscripts. AI assistance ranges from semi-automated grammar and citation tools to fully automated paper generation: the most commercially mature yet ethically contested stage.}
\vspace{2pt}
\resizebox{\linewidth}{!}{\begin{tabular}{c| >{\centering\arraybackslash}p{3cm} r r |c| >{\centering\arraybackslash}p{2.5cm} |c| p{7.2cm}}
\toprule
\rowcolor{tableheader!10}
\textbf{\#} & \textbf{Method} & \textbf{Ref} & \textbf{Venue} & \textbf{Link} & \textbf{Category} & \textbf{GitHub} & \textbf{Evaluation} \\
\midrule
\multicolumn{8}{@{}l}{\cellcolor{S5color!65}\textbf{\textsf{\textcolor{white}{~~Semi-Automated Writing Assistance}}}} \\
\addlinespace[1pt]
{1}  & CoAuthor & \cite{coauthor2022} & {\small arXiv'22} & \href{https://arxiv.org/abs/2201.06796}{\faExternalLink} & Collaborative    & -                                                                     & Human--AI collaborative writing workflows \\
{2}  & Script\&Shift & {\small \cite{siddiqui2025scriptshift}} & {\small CHI'25} & \href{https://arxiv.org/abs/2502.07440}{\faExternalLink}   & Source Transform & -                                                                     & CHI Honorable Mention; preserves cognitive engagement \\
{3}  & AI Writing Study & \cite{siddiqui2025aiwriting} & {\small AIED'25} & \href{https://arxiv.org/abs/2506.20595}{\faExternalLink} & Empirical Study  & -                                                                     & 90-student RCT; purposeful AI fosters writing \\
{4}  & OpenDraft & \cite{opendraft2025} & - & \href{https://github.com/federicodeponte/opendraft}{\faExternalLink} & Full Draft Gen.   & \githubicon{https://github.com/federicodeponte/opendraft}                & 19 agents; 20K+ words in 10 min; verified citations \\
{5}  & DraftMarks & \cite{siddiqui2025draftmarks} & {\small arXiv'25} & \href{https://arxiv.org/abs/2509.23505}{\faExternalLink} & Transparency     & -                                                                     & Skeuomorphic visual traces for AI process transparency \\
{6}  & PaperDebugger & \cite{paperdebugger2025} & {\small arXiv'25} & \href{https://arxiv.org/abs/2512.02589}{\faExternalLink} & In-Editor Assist & \githubicon{https://github.com/PaperDebugger/PaperDebugger}              & Multi-agent Overleaf plugin (Reviewer+Enhancer+Scorer) \\
{7}  & ScholarCopilot & \cite{scholarcop2025} & {\small arXiv'25} & \href{https://arxiv.org/abs/2504.00824}{\faExternalLink} & Citation Assist  & -                                                                     & 40.1\% top-1 citation accuracy (vs 15.0\% E5-Mistral) \\
{8}  & XtraGPT & \cite{xtragpt2025} & {\small arXiv'25} & \href{https://arxiv.org/abs/2505.11336}{\faExternalLink} & Post-Writing     & -                                                                     & 1.5B--14B models; 7K papers; 140K revision pairs \\
{9}  & LimAgents & \cite{limagents2026} & {\small arXiv'26} & \href{https://arxiv.org/abs/2601.11578}{\faExternalLink} & Limitations Gen. & -                                                                     & OpenReview comments + citation network integration \\
\addlinespace[3pt]
\multicolumn{8}{@{}l}{\cellcolor{S5color!65}\textbf{\textsf{\textcolor{white}{~~Fully Automated Paper Generation}}}} \\
\addlinespace[1pt]
{10} & CycleResearcher & \cite{cycleresearcher2024} & {\small ICLR'25} & \href{https://arxiv.org/abs/2411.00816}{\faExternalLink} & E2E Gen.   & -                                                                     & 5.36 ICLR scale (vs 5.24 preprint, 5.69 accepted) \\
{11} & Agent Laboratory & \cite{schmidgall2025agentlab} & {\small EMNLP'25} & \href{https://aclanthology.org/2025.findings-emnlp.320/}{\faExternalLink} & E2E Gen.   & -                                                                     & \$2--13/paper; 84\% cost reduction; 3.5--4.0 score \\
{12} & FutureGen & \cite{futuregen2025} & {\small arXiv'25} & \href{https://arxiv.org/abs/2503.16561}{\faExternalLink} & Section Gen.      & -                                                                     & RAG-based Future Work section generation \\
{13} & AI Scientist & \cite{lu2024aiscientist} & {\small Nature'26} & \href{https://arxiv.org/abs/2408.06292}{\faExternalLink} & E2E Gen.   & \githubicon{https://github.com/SakanaAI/AI-Scientist}                    & \$15/paper; end-to-end across 3 ML subfields \\
{14} & APRES & \cite{apres2026} & {\small arXiv'26} & \href{https://arxiv.org/abs/2603.03142}{\faExternalLink} & Rubric Revision  & -                                                                     & 79\% expert preference; citation-predictive rubrics \\
\addlinespace[3pt]
\multicolumn{8}{@{}l}{\cellcolor{S5color!65}\textbf{\textsf{\textcolor{white}{~~Societal Analysis}}}} \\
\addlinespace[1pt]
{15} & AI Writing Adoption & \cite{aicontractfocus2025} & {\small Nature'26} & \href{https://www.nature.com/articles/s41586-025-08681-8}{\faExternalLink} & Measurement      & -                                                                     & 41.3M papers; AI expands impact but contracts focus \\
{16} & Nature AI Survey & \cite{aireviewsurvey2025} & {\small Nature'26} & \href{https://www.nature.com/articles/d41586-025-04066-5}{\faExternalLink} & Survey           & -                                                                     & 57\% of researchers use AI in peer review \\
\addlinespace[3pt]
\multicolumn{8}{@{}l}{\cellcolor{S5color!65}\textbf{\textsf{\textcolor{white}{~~Writing Quality and AI Detection Assessment}}}} \\
\addlinespace[1pt]
{17} & Mapping LLM Use & \cite{liang2024mapping} & {\small arXiv'24} & \href{https://arxiv.org/abs/2404.01268}{\faExternalLink} & Detection        & -                                                                     & Up to 17.5\% of CS papers AI-modified \\
{18} & CycleReviewer & \cite{cycleresearcher2024} & {\small ICLR'25} & \href{https://arxiv.org/abs/2411.00816}{\faExternalLink} & AI Judge         & -                                                                     & 26.89\% MAE reduction vs individual human reviewers \\
{19} & Stanford Agentic & \cite{stanfordreviewer2025} & {\small Web'25} & \href{https://paperreview.ai/tech-overview}{\faExternalLink} & AI Judge         & -                                                                     & $\rho=0.42$ vs human $\rho=0.41$; matches consistency \\
{20} & SciIG & \cite{sciig2025} & {\small arXiv'25} & \href{https://arxiv.org/abs/2508.14273}{\faExternalLink} & Writing Bench    & -                                                                     & NAACL/ICLR 2025 introduction writing benchmark$^\dagger$ \\
{21} & Watermarking & \cite{watermarking2025} & {\small arXiv'25} & \href{https://arxiv.org/abs/2503.15772}{\faExternalLink} & Detection        & -                                                                     & Near-zero false-positive rate under controlled conditions \\
{22} & PaperWritingBench & \cite{paperwritingbench2026} & {\small arXiv'26} & \href{https://arxiv.org/abs/2604.05018}{\faExternalLink} & Benchmark        & -                                                                     & 200 reverse-engineered top-tier conference papers \\
\bottomrule
\end{tabular}}
\end{table*}
\endgroup



\clearpage
\subsection{Phase 3: Validation}

% ==================== Appendix Table: S6 Peer Review ====================
\begingroup
\renewcommand{\arraystretch}{1.15}
\setlength{\tabcolsep}{3pt}
\footnotesize
\rowcolors{2}{white}{S6color!6}
\begin{table*}[!ht]
\centering
\caption{\textbf{Comprehensive inventory: \Ssix Peer Review.} $^\dagger$Evaluation information uncertain.}
\vspace{-7pt}
\label{tab:appendix_s6}
\stagecard{\Ssix}{Peer Review}{S6color}{figures/icons/s6_review.png}{%
Automated Review Generation, reviewer--paper matching, and review quality assessment. AI systems can produce structured critiques and predict acceptance decisions, though leniency bias and adversarial vulnerabilities persist.}
\vspace{2pt}
\resizebox{\linewidth}{!}{\begin{tabular}{c| >{\centering\arraybackslash}p{3cm} r r |c| >{\centering\arraybackslash}p{2.1cm} |c| p{7.8cm}}
\toprule
\rowcolor{tableheader!10}
\textbf{\#} & \textbf{Method} & \textbf{Ref} & \textbf{Venue} & \textbf{Link} & \textbf{Category} & \textbf{GitHub} & \textbf{Evaluation} \\
\midrule
\multicolumn{8}{@{}l}{\cellcolor{S6color!65}\textbf{\textsf{\textcolor{white}{~~Automated Review Generation}}}} \\
\addlinespace[1pt]
{1}  & ChatReviewer & \cite{chatreviewer2023} & {\small GitHub'23} & \href{https://github.com/nishiwen1214/ChatReviewer}{\faExternalLink} & Review Gen.       & \githubicon{https://github.com/nishiwen1214/ChatReviewer}                & ChatGPT-based strengths/weaknesses analysis \\
{2}  & AI-Peer-Review & \cite{poldrack2024aireview} & {\small GitHub'24} & \href{https://github.com/poldrack/ai-peer-review}{\faExternalLink} & Review Gen.       & \githubicon{https://github.com/poldrack/ai-peer-review}                  & Multi-LLM independent reviews + meta-review synthesis \\
{3}  & MARG & \cite{darcy2024marg} & {\small arXiv'24} & \href{https://arxiv.org/abs/2401.04259}{\faExternalLink} & Review Gen.       & -                                                                     & 3.7 good comments/paper (2.2$\times$ over baseline) \\
{4}  & Reviewer2 & \cite{reviewer2024} & {\small arXiv'24} & \href{https://arxiv.org/abs/2402.10886}{\faExternalLink} & Review Gen.       & -                                                                     & Two-stage prompt-based review aspect modeling \\
{5}  & ReviewRL & \cite{reviewrl2025} & {\small EMNLP'25} & \href{https://aclanthology.org/2025.emnlp-main.857/}{\faExternalLink} & Review Gen.       & -                                                                     & RL + RAG; factually grounded reviews \\
{6}  & DeepReviewer & \cite{deepreviewer2025} & {\small arXiv'25} & \href{https://arxiv.org/abs/2503.08569}{\faExternalLink} & Review Gen.       & -                                                                     & 88.21\% win rate vs GPT-o1; 64\% accept/reject acc. \\
{7}  & OpenReviewer & \cite{openreviewer2025} & {\small NAACL'25} & \href{https://aclanthology.org/2025.naacl-demo.44/}{\faExternalLink} & Review Gen.       & -                                                                     & Llama-8B fine-tuned on 79K expert reviews \\
{8}  & REMOR & \cite{remor2025} & {\small arXiv'25} & \href{https://arxiv.org/abs/2505.11718}{\faExternalLink} & Review Gen.       & -                                                                     & Multi-objective RL review optimization \\
{9}  & ScholarPeer & \cite{scholarpeer2026} & {\small arXiv'26} & \href{https://arxiv.org/abs/2601.22638}{\faExternalLink} & Review Gen.       & -                                                                     & Context-aware multi-agent; literature verification \\
\addlinespace[3pt]
\multicolumn{8}{@{}l}{\cellcolor{S6color!65}\textbf{\textsf{\textcolor{white}{~~Meta-Review \& Reviewer Matching}}}} \\
\addlinespace[1pt]
{10} & AgentReview & \cite{agentreview2024} & {\small EMNLP'24} & \href{https://aclanthology.org/2024.emnlp-main.70/}{\faExternalLink} & Meta-Review      & -                                                                     & Full review lifecycle simulation; social/authority bias \\
{11} & Meta-Review LLMs & \cite{metareviewllm2024} & {\small NAACL'25} & \href{https://aclanthology.org/2025.naacl-long.395/}{\faExternalLink} & Meta-Review      & -                                                                     & 40 ICLR papers; GPT-3.5/LLaMA-2/PaLM-2 compared \\
{12} & RATE & \cite{rate2026} & {\small arXiv'26} & \href{https://arxiv.org/abs/2601.19637}{\faExternalLink} & Matching         & -                                                                     & Expertise-based matching via profile distillation \\
\addlinespace[3pt]
\multicolumn{8}{@{}l}{\cellcolor{S6color!65}\textbf{\textsf{\textcolor{white}{~~Adversarial Attacks \& Bias Analysis}}}} \\
\addlinespace[1pt]
{13} & Raina~\etal & \cite{raina2024adversarialjudge} & {\small EMNLP'24} & \href{https://arxiv.org/abs/2402.14016}{\faExternalLink} & Adversarial      & -                                                                     & Benign adjectives as universal adversarial triggers \\
{14} & AI Review Lottery & \cite{ailottery2024} & {\small arXiv'24} & \href{https://arxiv.org/abs/2405.02150}{\faExternalLink} & Bias Analysis    & -                                                                     & 15.8\% ICLR reviews AI-assisted; +4.9pp borderline \\
{15} & Ye~\etal & \cite{ye2024peerrisks} & {\small arXiv'24} & \href{https://arxiv.org/abs/2412.01708}{\faExternalLink} & Adversarial      & -                                                                     & Scores inflated to $\sim$8; 5\% manipulation flips 12\% \\
{16} & Breaking the Reviewer & \cite{breakingreviewer2025} & {\small arXiv'25} & \href{https://arxiv.org/abs/2506.11113}{\faExternalLink} & Adversarial     & -                                                                     & Systematic adversarial robustness evaluation \\
{17} & LLM Reviewer Bias & \cite{llmreviewer2025} & {\small arXiv'25} & \href{https://arxiv.org/abs/2509.09912}{\faExternalLink} & Bias Analysis    & -                                                                     & 1,441 papers; 95.8\% rejected misclassified as acceptable \\
{18} & Prompt Injection & \cite{promptinjectionreview2025} & {\small arXiv'25} & \href{https://arxiv.org/abs/2509.10248}{\faExternalLink} & Adversarial     & -                                                                     & White-text injection; up to 100\% acceptance scores \\
{19} & Sahoo~\etal & \cite{sahoo2025indirect} & {\small arXiv'25} & \href{https://arxiv.org/abs/2512.10449}{\faExternalLink} & Adversarial      & -                                                                     & +13.95 on Mistral; 13 LLMs; 15 attack strategies \\
{20} & Zhou~\etal & \cite{zhou2025positiveprompt} & {\small arXiv'25} & \href{https://arxiv.org/abs/2511.01287}{\faExternalLink} & Adversarial      & -                                                                     & +1.24 to +2.80 from hype prose; 10.00 under iteration \\
\addlinespace[3pt]
\multicolumn{8}{@{}l}{\cellcolor{S6color!65}\textbf{\textsf{\textcolor{white}{~~Detection \& Policy}}}} \\
\addlinespace[1pt]
{21} & AI Detection & \cite{aidetectionreview2025} & {\small arXiv'25} & \href{https://arxiv.org/abs/2502.19614}{\faExternalLink} & Detection        & -                                                                     & 788,984 AI-written reviews; 18 detection algorithms \\
{22} & AI Use Rejects & \cite{aiuserejects2026} & {\small Nature'26} & \href{https://www.nature.com/articles/d41586-026-00893-2}{\faExternalLink} & Policy           & -                                                                     & 497 papers rejected ($\sim$2\% of submissions) \\
{23} & Nature AI Survey & \cite{aireviewsurvey2025} & {\small Nature'26} & \href{https://www.nature.com/articles/d41586-025-04066-5}{\faExternalLink} & Survey           & -                                                                     & 1,600 academics; 57\% use AI in peer review \\
{24} & Policy Enforcement & \cite{reviewpolicyenforce2026} & {\small arXiv'26} & \href{https://arxiv.org/abs/2603.20450}{\faExternalLink} & Policy          & -                                                                     & All 5 SOTA detectors misclassify LLM-polished reviews \\
{25} & Reviewer Feedback & \cite{reviewerfeedback2026} & {\small CHI'26} & \href{https://doi.org/10.1145/3772318.3791431}{\faExternalLink} & Empirical        & -                                                                     & ICLR 2025 live process; reviewer engagement study$^\dagger$ \\
\addlinespace[3pt]
\multicolumn{8}{@{}l}{\cellcolor{S6color!65}\textbf{\textsf{\textcolor{white}{~~Review Consistency and Bias Assessment}}}} \\
\addlinespace[1pt]
{26} & Review Survey & \cite{zhuang2025asprsurvey} & {\small IF'25} & \href{https://www.sciencedirect.com/science/article/pii/S1566253525004051}{\faExternalLink} & Survey     & -                                                                     & Comprehensive taxonomy of review methods \\
{27} & Stanford Agentic & \cite{stanfordreviewer2025} & {\small Web'25} & \href{https://paperreview.ai/tech-overview}{\faExternalLink} & Quality          & -                                                                     & $\rho=0.42$ vs human $\rho=0.41$; matches consistency \\
{28} & ClaimCheck & \cite{claimcheck2025} & {\small EMNLP'25} & \href{https://aclanthology.org/2025.findings-emnlp.1185/}{\faExternalLink} & Quality          & -                                                                     & LLM critique grounding; gaps in factual basis \\
{29} & ReViewGraph & \cite{reviewgraph2025} & {\small AAAI'26} & \href{https://arxiv.org/abs/2511.08317}{\faExternalLink} & Quality          & -                                                                     & +15.73\% avg improvement via heterogeneous graph \\
{30} & ReviewAgents & \cite{reviewagents2025} & {\small arXiv'25} & \href{https://arxiv.org/abs/2503.08506}{\faExternalLink} & Quality          & -                                                                     & 37,403 papers; 142,324 reviews; Review-CoT dataset \\
{31} & ICLR 2025 Study & \cite{iclr2025reviewstudy} & {\small NMI'26} & \href{https://www.nature.com/articles/s42256-026-01188-x}{\faExternalLink} & Quality       & -                                                                     & 22,467 reviews; 89\% quality improved; no acceptance effect \\
\bottomrule
\end{tabular}}
\end{table*}
\endgroup


\clearpage
% ==================== Appendix Table: S7 Rebuttal ====================
\begingroup
\renewcommand{\arraystretch}{1.15}
\setlength{\tabcolsep}{3pt}
\footnotesize
\rowcolors{2}{white}{S7color!6}
\begin{table*}[!ht]
\centering
\caption{\textbf{Comprehensive inventory: \Sseven Rebuttal \& Revision.} $^\dagger$Evaluation information uncertain.}
\vspace{-7pt}
\label{tab:appendix_s7}
\stagecard{\Sseven}{Rebuttal \& Revision}{S7color}{figures/icons/s7_rebuttal.png}{%
Analyzing reviewer comments and generating evidence-grounded responses. The newest frontier in AI auto-research, rebuttal automation is decisive for roughly one in five borderline submissions at major venues.}
\vspace{2pt}
\resizebox{\linewidth}{!}{\begin{tabular}{c| >{\centering\arraybackslash}p{3cm} r r |c| >{\centering\arraybackslash}p{2.4cm} |c| p{7.5cm}}
\toprule
\rowcolor{tableheader!10}
\textbf{\#} & \textbf{Method} & \textbf{Ref} & \textbf{Venue} & \textbf{Link} & \textbf{Category} & \textbf{GitHub} & \textbf{Evaluation} \\
\midrule
\multicolumn{8}{@{}l}{\cellcolor{S7color!65}\textbf{\textsf{\textcolor{white}{~~Reviewer Comment Analysis}}}} \\
\addlinespace[1pt]
{1}  & ReviewMT & \cite{reviewmt2024} & {\small arXiv'24} & \href{https://arxiv.org/abs/2406.05688}{\faExternalLink} & Analysis         & -                                                                     & 26,841 papers; 92,017 reviews; multi-turn dialogue \\
{2}  & ICLR Rebuttal Study & \cite{iclr_rebuttal2025} & {\small arXiv'25} & \href{https://arxiv.org/abs/2511.15462}{\faExternalLink} & Analysis         & -                                                                     & ICLR 2024--2025; score transition analysis \\
\addlinespace[3pt]
\multicolumn{8}{@{}l}{\cellcolor{S7color!65}\textbf{\textsf{\textcolor{white}{~~Automated Rebuttal Generation}}}} \\
\addlinespace[1pt]
{3}  & ReviewerToo & \cite{reviewertoo2025} & {\small arXiv'25} & \href{https://arxiv.org/abs/2510.08867}{\faExternalLink} & Modular Pipeline & -                                                                     & 81.8\% accept/reject accuracy (vs 83.9\% human) \\
{4}  & RebuttalAgent & \cite{rebuttalagent2026} & {\small ICLR'26} & \href{https://arxiv.org/abs/2601.15715}{\faExternalLink} & Rebuttal Gen.     & \githubicon{https://github.com/Zhitao-He/RebuttalAgent}             & 18.3\% avg improvement; ToM-grounded \\
{5}  & Author-in-the-Loop & \cite{ruan2026authorinloop} & {\small ACL'26} & \href{https://arxiv.org/abs/2602.11173}{\faExternalLink} & Author-Aware     & -                                                                     & Integrates author expertise and intent \\
{6}  & DRPG & \cite{drpg2026} & {\small arXiv'26} & \href{https://arxiv.org/abs/2601.18081}{\faExternalLink} & Rebuttal Gen.     & \githubicon{https://github.com/ulab-uiuc/DRPG-RebuttalAgent}             & 98\%+ planning accuracy; surpasses avg human quality \\
{7}  & Paper2Rebuttal & \cite{paper2rebuttal2026} & {\small arXiv'26} & \href{https://arxiv.org/abs/2601.14171}{\faExternalLink} & Rebuttal Gen.     & -                                                                     & Evidence-centric rebuttal planning \\
\addlinespace[3pt]
\multicolumn{8}{@{}l}{\cellcolor{S7color!65}\textbf{\textsf{\textcolor{white}{~~Rebuttal Effectiveness Assessment}}}} \\
\addlinespace[1pt]
{8}  & Re$^2$ & \cite{re2dataset2025} & {\small arXiv'25} & \href{https://arxiv.org/abs/2505.07920}{\faExternalLink} & Dataset          & -                                                                     & 19,926 submissions; 70,668 reviews; 53,818 rebuttals \\
{9}  & Commitment Checklist & \cite{rebuttalcommitment2026} & {\small arXiv'26} & \href{https://arxiv.org/abs/2603.00003}{\faExternalLink} & Benchmark     & -                                                                     & 11.8 commitments/paper; $\sim$25\% unfulfilled \\
{10} & Re$^3$Align & \cite{ruan2026authorinloop} & {\small ACL'26} & \href{https://arxiv.org/abs/2602.11173}{\faExternalLink} & Dataset     & -                                                                     & First large-scale aligned review--response--revision triplets \\
\bottomrule
\end{tabular}}
\end{table*}
\endgroup



\clearpage
\subsection{Phase 4: Dissemination}

% ==================== Appendix Table: S8 Dissemination ====================
\begingroup
\renewcommand{\arraystretch}{1.15}
\setlength{\tabcolsep}{3pt}
\footnotesize
\rowcolors{2}{white}{S8color!6}
\begin{table*}[!ht]
\centering
\caption{\textbf{Comprehensive inventory: \Seight Dissemination (Paper2X).} $^\dagger$Evaluation information uncertain.}
\vspace{-7pt}
\label{tab:appendix_s8}
\stagecard{\Seight}{Dissemination / Paper2X}{S8color}{figures/icons/s8_dissemination.png}{%
Converting papers into posters, slides, videos, websites, and social media content. Each output format targets a different audience and demands its own design logic, AI tool chain, and trust considerations.}
\vspace{2pt}
\resizebox{\linewidth}{!}{\begin{tabular}{c| >{\centering\arraybackslash}p{3cm} r r |c| >{\centering\arraybackslash}p{2.2cm} |c| p{7.2cm}}
\toprule
\rowcolor{tableheader!10}
\textbf{\#} & \textbf{Method} & \textbf{Ref} & \textbf{Venue} & \textbf{Link} & \textbf{Category} & \textbf{GitHub} & \textbf{Evaluation} \\
\midrule
\multicolumn{8}{@{}l}{\cellcolor{S8color!65}\textbf{\textsf{\textcolor{white}{~~Paper2Poster}}}} \\
\addlinespace[1pt]
{1}  & P2P & \cite{p2p2025} & {\small ICLR'26} & \href{https://openreview.net/forum?id=JojyT9niJL}{\faExternalLink} & Paper2Poster     & -                                                                     & P2PInstruct 30K+ examples; 3 specialized agents \\
{2}  & Paper2Poster & \cite{paper2poster2025} & {\small NeurIPS'25} & \href{https://openreview.net/forum?id=p0E74lpRBD}{\faExternalLink} & Paper2Poster     & \githubicon{https://github.com/Paper2Poster/Paper2Poster}                & \$0.005/poster; 87\% fewer tokens vs GPT-4o \\
{3}  & PosterForest & \cite{choi2025posterforest} & {\small arXiv'25} & \href{https://arxiv.org/abs/2508.21720}{\faExternalLink} & Paper2Poster     & -                                                                     & Hierarchical multi-agent collaboration$^\dagger$ \\
{4}  & PosterGen & \cite{postergen2025} & {\small arXiv'25} & \href{https://arxiv.org/abs/2508.17188}{\faExternalLink} & Paper2Poster     & -                                                                     & Aesthetic-aware multi-agent generation$^\dagger$ \\
{5}  & APEX & \cite{apex2026} & {\small arXiv'26} & \href{https://arxiv.org/abs/2601.04794}{\faExternalLink} & Paper2Poster     & \githubicon{https://github.com/Breesiu/APEX}                             & First agentic interactive poster editing \\
{6}  & PosterOmni & \cite{posteromni2026} & {\small arXiv'26} & \href{https://arxiv.org/abs/2602.12127}{\faExternalLink} & Paper2Poster     & -                                                                     & 6 unified poster tasks; outperforms open-source \\
\addlinespace[3pt]
\multicolumn{8}{@{}l}{\cellcolor{S8color!65}\textbf{\textsf{\textcolor{white}{~~Paper2Slides}}}} \\
\addlinespace[1pt]
{7}  & DOC2PPT & \cite{doc2ppt2022} & {\small AAAI'22} & \href{https://ojs.aaai.org/index.php/AAAI/article/view/19943}{\faExternalLink} & Paper2Slides     & -                                                                     & 5,873 paired document--slide decks \\
{8}  & PPTAgent & \cite{pptagent2025} & {\small EMNLP'25} & \href{https://arxiv.org/abs/2501.03936}{\faExternalLink} & Paper2Slides     & \githubicon{https://github.com/icip-cas/PPTAgent}                        & PPTEval benchmark; 10,448 curated presentations \\
{9}  & AutoPresent & \cite{autopresent2025} & {\small CVPR'25} & \href{https://arxiv.org/abs/2501.00912}{\faExternalLink} & Paper2Slides     & -                                                                     & 8B Llama model; SlidesBench (7K train, 585 test) \\
{10} & Paper2Slides & \cite{paper2slides2025} & {\small GitHub'25} & \href{https://github.com/HKUDS/Paper2Slides}{\faExternalLink} & Paper2Slides     & \githubicon{https://github.com/HKUDS/Paper2Slides}                       & 4-stage RAG pipeline; one-click conversion \\
{11} & Auto-Slides & \cite{autoslides2025} & {\small arXiv'25} & \href{https://arxiv.org/abs/2509.11062}{\faExternalLink} & Paper2Slides     & -                                                                     & Multi-agent Beamer generation; interactive editing$^\dagger$ \\
{12} & PASS & \cite{pass2025} & {\small arXiv'25} & \href{https://arxiv.org/abs/2501.06497}{\faExternalLink} & Paper2Slides     & -                                                                     & First combined slides + AI audio delivery$^\dagger$ \\
{13} & SlideGen & \cite{slidegen2025} & {\small arXiv'25} & \href{https://arxiv.org/abs/2512.04529}{\faExternalLink} & Paper2Slides     & -                                                                     & Multi-agent VLM coordination; editable PPTX output \\
{14} & Talk to Your Slides & \cite{talkslides2025} & {\small arXiv'25} & \href{https://arxiv.org/abs/2505.11604}{\faExternalLink} & Paper2Slides     & -                                                                     & +34\% instruction fidelity; 87\% lower cost vs GUI \\
{15} & SlideTailor & \cite{slidetailor2025} & {\small AAAI'26} & \href{https://arxiv.org/abs/2512.20292}{\faExternalLink} & Paper2Slides     & \githubicon{https://github.com/nusnlp/SlideTailor}                       & User-preference conditioned; chain-of-speech \\
{16} & DeepPresenter & \cite{deeppresenter2026} & {\small arXiv'26} & \href{https://arxiv.org/abs/2602.22839}{\faExternalLink} & Paper2Slides     & \githubicon{https://github.com/icip-cas/PPTAgent}                        & 9B model competitive with frontier at lower cost \\
{17} & Office Raccoon & \cite{sensetime2026} & {\small Web'26} & \href{https://www.sensetime.com/en/news-detail/51170569}{\faExternalLink} & Paper2Slides     & -                                                                     & Page-level editing; template/brand-guideline learning$^\dagger$ \\
\addlinespace[3pt]
\multicolumn{8}{@{}l}{\cellcolor{S8color!65}\textbf{\textsf{\textcolor{white}{~~Paper2Video}}}} \\
\addlinespace[1pt]
{18} & Preacher & \cite{preacher2025} & {\small ICCV'25} & \href{https://github.com/Gen-Verse/Paper2Video}{\faExternalLink} & Paper2Video      & \githubicon{https://github.com/Gen-Verse/Paper2Video}                    & Top-down decomposition; 5 research fields \\
{19} & Paper2Video & \cite{paper2video2025} & {\small arXiv'25} & \href{https://arxiv.org/abs/2510.05096}{\faExternalLink} & Paper2Video      & \githubicon{https://github.com/showlab/Paper2Video}                      & 101 paper--video pairs; +10\% PresentQuiz accuracy \\
{20} & PresentAgent & \cite{presentagent2025} & {\small EMNLP'25} & \href{https://aclanthology.org/2025.emnlp-demos.58/}{\faExternalLink} & Paper2Video      & \githubicon{https://github.com/AIGeeksGroup/PresentAgent}                & PresentEval benchmark; approaches human-level \\
\addlinespace[3pt]
\multicolumn{8}{@{}l}{\cellcolor{S8color!65}\textbf{\textsf{\textcolor{white}{~~Paper2Web \& Social Media}}}} \\
\addlinespace[1pt]
{21} & Paper2Web & \cite{paper2web2025} & {\small arXiv'25} & \href{https://arxiv.org/abs/2510.15842}{\faExternalLink} & Paper2Web        & \githubicon{https://github.com/YuhangChen1/Paper2All}                    & 10,716 papers; multimedia-rich academic homepages \\
\addlinespace[3pt]
\multicolumn{8}{@{}l}{\cellcolor{S8color!65}\textbf{\textsf{\textcolor{white}{~~Fidelity and Adoption Assessment}}}} \\
\addlinespace[1pt]
{22} & PPTEval & \cite{pptagent2025} & {\small EMNLP'25} & \href{https://arxiv.org/abs/2501.03936}{\faExternalLink} & Benchmark        & \githubicon{https://github.com/icip-cas/PPTAgent}                        & Content, design, coherence; 10,448 presentations \\
{23} & PresentQuiz & \cite{paper2video2025} & {\small arXiv'25} & \href{https://arxiv.org/abs/2510.05096}{\faExternalLink} & Benchmark        & \githubicon{https://github.com/showlab/Paper2Video}                      & 101 paper--video pairs; +10\% over human on comprehension \\
{24} & PresentEval & \cite{presentagent2025} & {\small EMNLP'25} & \href{https://aclanthology.org/2025.emnlp-demos.58/}{\faExternalLink} & Benchmark        & \githubicon{https://github.com/AIGeeksGroup/PresentAgent}                & End-to-end narrated video quality; near human-level \\
\bottomrule
\end{tabular}}
\end{table*}
\endgroup



\clearpage
\subsection{Cross-Phase: End-to-End Systems}

% ==================== Appendix Table: Cross-Phase End-to-End Systems ====================
\begingroup
\renewcommand{\arraystretch}{1.15}
\setlength{\tabcolsep}{3pt}
\footnotesize
\rowcolors{2}{white}{tableheader!6}
\begin{table*}[!ht]
\centering
\caption{\textbf{Comprehensive inventory: End-to-End and Cross-Phase Systems.} Systems that span multiple stages of the research lifecycle. $^\dagger$Evaluation information uncertain.}
\vspace{-7pt}
\label{tab:appendix_e2e}
\vspace{2pt}
\resizebox{\linewidth}{!}{\begin{tabular}{c| >{\centering\arraybackslash}p{3cm} r r |c| >{\centering\arraybackslash}p{2.4cm} |c| p{7.2cm}}
\toprule
\rowcolor{tableheader!10}
\textbf{\#} & \textbf{Method} & \textbf{Ref} & \textbf{Venue} & \textbf{Link} & \textbf{Category} & \textbf{GitHub} & \textbf{Evaluation} \\
\midrule
\multicolumn{8}{@{}l}{\cellcolor{tableheader!65}\textbf{\textsf{\textcolor{white}{~~Fully Automated Research Systems}}}} \\
\addlinespace[1pt]
{1} & AI Scientist & \cite{lu2024aiscientist} & {\small arXiv'24} & \href{https://arxiv.org/abs/2408.06292}{\faExternalLink} & E2E Pipeline & \githubicon{https://github.com/SakanaAI/AI-Scientist} & Pioneered E2E at \$15/paper; ICLR-scale review \\
{2} & Agent Laboratory & \cite{schmidgall2025agentlab} & {\small EMNLP'25} & \href{https://aclanthology.org/2025.findings-emnlp.320/}{\faExternalLink} & E2E Pipeline & - & \$2--13/paper; 3.5--4.0 NeurIPS scale \\
{3} & AI-Researcher & \cite{airesearcher2025} & {\small arXiv'25} & \href{https://arxiv.org/abs/2505.18705}{\faExternalLink} & E2E Pipeline & - & Scientist-Bench; approaches human-level quality \\
{4} & CycleResearcher & \cite{cycleresearcher2024} & {\small ICLR'25} & \href{https://arxiv.org/abs/2411.00816}{\faExternalLink} & E2E Pipeline & - & 5.36 ICLR scale; cyclic write--review \\
{5} & AI Scientist v2 & \cite{yamada2025aiscientistv2} & {\small arXiv'25} & \href{https://arxiv.org/abs/2504.08066}{\faExternalLink} & E2E + Tree Search & \githubicon{https://github.com/SakanaAI/AI-Scientist-v2} & ICLR 2025 ICBINB workshop; score 6.33 \\
{6} & Kosmos & \cite{kosmos2025} & {\small arXiv'25} & \href{https://arxiv.org/abs/2511.02824}{\faExternalLink} & E2E Pipeline & - & 79.4\% claim accuracy; 7 discoveries; 4 domains \\
{7} & Dolphin & \cite{dolphin2025} & {\small ACL'25} & \href{https://aclanthology.org/2025.acl-long.1056/}{\faExternalLink} & E2E Pipeline & - & Closed-loop auto-research pipeline \\
{8} & CodeScientist & \cite{codescientist2025} & {\small ACL'25} & \href{https://aclanthology.org/2025.findings-acl.692/}{\faExternalLink} & E2E Pipeline & \githubicon{https://github.com/allenai/codescientist} & Hypothesis to verification; closed-loop \\
{9} & InternAgent & \cite{novelseek2025} & {\small arXiv'25} & \href{https://arxiv.org/abs/2505.16938}{\faExternalLink} & E2E Pipeline & \githubicon{https://github.com/InternScience/InternAgent} & Closed-loop hypothesis to verification \\
{10} & freephdlabor & \cite{freephdlabor2025} & {\small arXiv'25} & \href{https://arxiv.org/abs/2510.15624}{\faExternalLink} & Multi-Agent & - & Personalized research group; continual automation \\
{11} & SciMaster & \cite{scimaster2025} & {\small arXiv'25} & \href{https://arxiv.org/abs/2507.05241}{\faExternalLink} & Multi-Agent & \githubicon{https://github.com/sjtu-sai-agents/X-Master} & General-purpose scientific AI agents \\
{12} & ARIS & \cite{aris2025} & {\small GitHub'26} & \href{https://github.com/wanshuiyin/Auto-claude-code-research-in-sleep}{\faExternalLink} & Skill Library & \githubicon{https://github.com/wanshuiyin/Auto-claude-code-research-in-sleep} & 31 skills; score 5.0$\to$7.5; 20+ GPU experiments \\
{13} & EvoScientist & \cite{evoscientist2026} & {\small arXiv'26} & \href{https://arxiv.org/abs/2603.08127}{\faExternalLink} & Multi-Agent & \githubicon{https://github.com/EvoScientist/EvoScientist} & 6 papers accepted at ICAIS'25 \\
{14} & UniScientist & \cite{uniscientist2026} & {\small Web'26} & \href{https://unipat.ai/blog/UniScientist}{\faExternalLink} & Multi-Agent & - & 30B open-source; beats Claude Opus 4.5 \\
{15} & ASI-Evolve & \cite{asievolve2026} & {\small GitHub'26} & \href{https://github.com/GAIR-NLP/ASI-Evolve}{\faExternalLink} & Multi-Agent & \githubicon{https://github.com/GAIR-NLP/ASI-Evolve} & +0.97 DeltaNet; +18 MMLU; +12.5 GRPO \\
{16} & AiScientist-LH & \cite{aiscientistlonghorizon2026} & {\small arXiv'26} & \href{https://arxiv.org/abs/2604.13018}{\faExternalLink} & E2E + Hierarchical & - & Long-horizon ML engineering; File-as-Bus \\
{17} & AutoSOTA & \cite{autosota2026} & {\small arXiv'26} & \href{https://arxiv.org/abs/2604.05550}{\faExternalLink} & E2E Pipeline & \githubicon{https://github.com/tsinghua-fib-lab/AutoSOTA} & Paper-to-code-to-SOTA optimization \\
{18} & CORAL & \cite{coral2026} & {\small arXiv'26} & \href{https://arxiv.org/abs/2604.01658}{\faExternalLink} & Multi-Agent & \githubicon{https://github.com/Human-Agent-Society/CORAL} & Multi-agent evolution; SOTA on 10 tasks \\
{19} & FARS & \cite{fars2026} & {\small Web'26} & \href{https://analemma.ai/blog/introducing-fars}{\faExternalLink} & Multi-Agent & - & 100 papers in 228h; avg 5.05 ICLR scale \\
{20} & AutoResearchClaw & \cite{autoresearchclaw2026} & {\small GitHub'26} & \href{https://github.com/aiming-lab/AutoResearchClaw}{\faExternalLink} & Multi-Agent & \githubicon{https://github.com/aiming-lab/AutoResearchClaw} & fully autonomous 23-stage pipeline \\
\addlinespace[3pt]
\multicolumn{8}{@{}l}{\cellcolor{tableheader!65}\textbf{\textsf{\textcolor{white}{~~Domain-Specific Systems}}}} \\
\addlinespace[1pt]
{21} & Coscientist & \cite{boiko2023coscientist} & {\small Nature'23} & \href{https://www.nature.com/articles/s41586-023-06792-0}{\faExternalLink} & Chemistry & - & Autonomous chemistry; LLM-driven tool use \\
{22} & AlphaFold~3 & \cite{abramson2024alphafold3} & {\small Nature'24} & \href{https://www.nature.com/articles/s41586-024-07487-w}{\faExternalLink} & Biology & - & Biomolecular structure prediction \\
{23} & ChemCrow & \cite{bran2024chemcrow} & {\small NMI'24} & \href{https://www.nature.com/articles/s42256-024-00832-8}{\faExternalLink} & Chemistry & - & Chemistry tool orchestration \\
{24} & Medical AI Scientist & \cite{medicalaiscientist2026} & {\small arXiv'26} & \href{https://arxiv.org/abs/2603.28589}{\faExternalLink} & Medicine & - & Clinical research automation \\
\addlinespace[3pt]
\multicolumn{8}{@{}l}{\cellcolor{tableheader!65}\textbf{\textsf{\textcolor{white}{~~Evolutionary \& Self-Improving Systems}}}} \\
\addlinespace[1pt]
{25} & ShinkaEvolve & \cite{shinkaevolve2025} & {\small arXiv'25} & \href{https://arxiv.org/abs/2509.19349}{\faExternalLink} & Evolutionary & \githubicon{https://github.com/SakanaAI/ShinkaEvolve} & Open-ended sample-efficient program evolution \\
{26} & Darwin Godel Machine & \cite{darwingodelmachine2025} & {\small arXiv'25} & \href{https://arxiv.org/abs/2505.22954}{\faExternalLink} & Self-Improving & \githubicon{https://github.com/jennyzzt/dgm} & Open-ended evolution of self-improving agents \\
\addlinespace[3pt]
\multicolumn{8}{@{}l}{\cellcolor{tableheader!65}\textbf{\textsf{\textcolor{white}{~~Research Platforms \& Infrastructure}}}} \\
\addlinespace[1pt]
{27} & R\&D-Agent & \cite{chen2025rdagent} & {\small arXiv'25} & \href{https://arxiv.org/abs/2505.14738}{\faExternalLink} & Infrastructure & \githubicon{https://github.com/microsoft/RD-Agent} & Researcher+Developer dual-agent; MLE-Bench top \\
{28} & autoresearch & \cite{karpathy2025autoresearch} & {\small GitHub'25} & \href{https://github.com/karpathy/autoresearch}{\faExternalLink} & Infrastructure & \githubicon{https://github.com/karpathy/autoresearch} & ${\sim}$12 exp/hour overnight \\
{29} & Google AI Co-scientist & \cite{gottweis2025aicoscientist} & {\small arXiv'25} & \href{https://arxiv.org/abs/2502.18864}{\faExternalLink} & Platform & - & Multi-agent hypothesis gen + validation \\
{30} & ResearchTown & \cite{researchtown2025} & {\small ICML'25} & \href{https://arxiv.org/abs/2412.17767}{\faExternalLink} & Multi-Agent & \githubicon{https://github.com/ulab-uiuc/research-town} & Simulates research community with LLM agents \\
{31} & AgentRxiv & \cite{schmidgall2025agentrxiv} & {\small arXiv'25} & \href{https://arxiv.org/abs/2503.18102}{\faExternalLink} & Multi-Agent & - & 11.4\% improvement on MATH-500 \\
{32} & LabClaw & \cite{labclaw2026} & {\small Web'26} & \href{https://labclaw-ai.github.io}{\faExternalLink} & Skill Library & \githubicon{https://github.com/wu-yc/LabClaw} & 206 biomedical skills; always-on autonomous lab agent \\
{33} & PiFlow & \cite{piflow2025} & {\small arXiv'25} & \href{https://arxiv.org/abs/2505.15047}{\faExternalLink} & Multi-Agent & - & Principle-aware scientific discovery \\
\addlinespace[3pt]
\bottomrule
\end{tabular}}
\vspace{-1cm}
\end{table*}
\endgroup




\clearpage\clearpage
\clearpage
\section{Survey Coverage Comparison \& Taxonomy Analysis}
\label{sec:appendix_surveys}

\Cref{tab:survey_comparison} compares our coverage with five closely related concurrent efforts. The goal is not to rank surveys, but to clarify how our lifecycle framework differs in scope and organization.

Our eight-stage framework subsumes several prior taxonomies while making two distinctions explicit: AI auto-research should be analyzed across the complete lifecycle, and stages should be grouped by epistemological function rather than only by task name or autonomy level.


\begin{wraptable}{r}{0.59\textwidth}
\vspace{-12pt}
\centering
\caption{%
  Comparison of survey coverage across our four-phase research lifecycle framework.
  \cmarkc~= in-depth coverage, \pmarkc~= partial coverage, \xmarkc~= not covered,
  {\textcolor{violet}{\scriptsize$\bigstar$\,new}}~= newly introduced stage.
}
\vspace{-6pt}
\label{tab:survey_comparison}
\renewcommand{\arraystretch}{1.15}
\setlength{\tabcolsep}{3pt}
\resizebox{0.56\textwidth}{!}{
\begin{tabular}{@{}l!{\vrule width 0.8pt}c!{\vrule width 0.8pt}ccccc@{}}
\toprule
\rowcolor{tableheader!10}
\textbf{Stage} &
\rotatebox{75}{\makecell[l]{\textbf{Ours}}} &
\rotatebox{75}{\makecell[l]{LLM4SR\\{\tiny 2501}}} &
\rotatebox{75}{\makecell[l]{PR Survey\\{\tiny 2501}}} &
\rotatebox{75}{\makecell[l]{Auto$\to$Auton\\{\tiny 2505}}} &
\rotatebox{75}{\makecell[l]{AI4Research\\{\tiny 2507}}} &
\rotatebox{75}{\makecell[l]{AI Scientists\\{\tiny 2510}}} \\
%% ---- Phase 1: Creation ----
\multicolumn{7}{l}{\cellcolor{S1color!8}\textbf{Phase~1: Creation}} \\
\makecell[tl]{\Sone: Idea Generation\\\textcolor{gray}{\scriptsize novelty, feasibility, multi-agent}}
  & \cmarkc & \cmarkc & \xmarkc & \cmarkc & \cmarkc & \cmarkc \\
\makecell[tl]{\Stwo: Literature Review\\\textcolor{gray}{\scriptsize retrieval, survey gen, deep research}}
  & \cmarkc & \xmarkc & \xmarkc & \cmarkc & \cmarkc & \cmarkc \\
\makecell[tl]{\Sthree: Coding \& Experiments\\\textcolor{gray}{\scriptsize paper-to-code, execution, analysis}}
  & \cmarkc & \cmarkc & \xmarkc & \cmarkc & \cmarkc & \cmarkc \\
\makecell[tl]{\Sfour: Figures \& Tables\\\textcolor{gray}{\scriptsize diagrams, plots, formulas}}
  & \cmarkc & \xmarkc & \xmarkc & \pmarkc & \pmarkc & \xmarkc \\
\midrule
%% ---- Phase 2: Writing ----
\multicolumn{7}{l}{\cellcolor{S5color!8}\textbf{Phase~2: Writing}} \\
\makecell[tl]{\Sfive: Paper Writing\\\textcolor{gray}{\scriptsize semi-auto, full-auto, detection}}
  & \cmarkc & \cmarkc & \xmarkc & \pmarkc & \cmarkc & \cmarkc \\
\midrule
%% ---- Phase 3: Validation ----
\multicolumn{7}{l}{\cellcolor{S6color!8}\textbf{Phase~3: Validation}} \\
\makecell[tl]{\Ssix: Peer Review\\\textcolor{gray}{\scriptsize auto-review, matching, quality}}
  & \cmarkc & \cmarkc & \cmarkc & \pmarkc & \cmarkc & \xmarkc \\
\makecell[tl]{\Sseven: Rebuttal \& Revision\\\textcolor{gray}{\scriptsize comment analysis, rebuttal gen}}
  & {\textcolor{violet}{\scriptsize$\bigstar$\,new}} & \xmarkc & \xmarkc & \xmarkc & \xmarkc & \xmarkc \\
\midrule
%% ---- Phase 4: Dissemination ----
\multicolumn{7}{l}{\cellcolor{S8color!8}\textbf{Phase~4: Dissemination}} \\
\makecell[tl]{\Seight: Dissemination\\\textcolor{gray}{\scriptsize poster, slides, video, social}}
  & {\textcolor{violet}{\scriptsize$\bigstar$\,new}} & \xmarkc & \xmarkc & \xmarkc & \xmarkc & \xmarkc \\
\midrule
\rowcolor{tableheader!6}
\textbf{Stages covered}
  & \textbf{8/8} & \textbf{4} & \textbf{1} & \textbf{5} & \textbf{5} & \textbf{4} \\
\bottomrule
\end{tabular}}
\vspace{-9pt}
\end{wraptable}
\begin{itemize}[leftmargin=*]
  \item \textbf{AI4Research}~\cite{ai4research2025} defines five task categories: Comprehension, Survey, Discovery, Writing, and Review. These overlap with our \Sone--\Sthree, \Sfive, and \Ssix. Our framework newly elevates \Sfour (Tables \& Figures), \Sseven (Rebuttal \& Revision), and \Seight (Dissemination) as independent lifecycle stages.

  \item \textbf{From Automation to Autonomy}~\cite{zheng2025automation} organizes systems by autonomy level, from tool-like assistance to scientist-level automation. This axis is complementary: each of our stages can be instantiated at different autonomy levels, while our framework specifies \emph{where} in the research lifecycle the system operates.

  \item \textbf{LLM4SR}~\cite{luo2025llm4sr} proposes a four-part view centered on hypothesis, experiment, writing, and review. This structure is close to ours, but does not separately model rebuttal and revision as a feedback stage. Our Validation phase separates \Ssix from \Sseven, making the review--response loop explicit.

  \item \textbf{Automated Scholarly Paper Review}~\cite{zhuang2025asprsurvey} provide in-depth coverage of review generation, quality assessment, and reviewer--paper matching. They are complementary to our work: they focus on \Ssix, while our framework places peer review within the broader lifecycle.

  \item \textbf{AI Scientist Survey}~\cite{tie2025survey} focus on autonomous or semi-autonomous scientific discovery, overlapping mainly with \Sone--\Sthree and \Sfive. Our framework extends this view by also covering scientific visualization, peer validation, rebuttal, and dissemination.
\end{itemize}

These comparisons show that prior taxonomies often list research tasks sequentially, while leaving functional distinctions and feedback loops implicit. Our four-phase framework makes these dependencies explicit. For example, \Ssix (\emph{Peer Review}) and \Sseven (\emph{Rebuttal \& Revision}) do not simply follow paper writing as isolated downstream steps; they can redirect the workflow back to \Sthree for additional experiments, \Sfour for revised figures or tables, and \Sfive for manuscript restructuring. Similarly, dissemination artifacts in \Seight may expose ambiguities in the original framing, requiring revisions to claims, explanations, or visual evidence. 

These cross-stage dependencies are central to real research practice and are especially important for AI-assisted workflows, where errors can propagate from generated ideas to experiments, from experiments to claims, and from claims to public-facing summaries. By organizing the field into \emph{Creation}, \emph{Writing}, \emph{Validation}, and \emph{Dissemination}, our framework highlights not only which stages are covered by existing systems, but also where evidence, claims, critique, and communication must remain aligned.